\documentclass[12pt, a4paper]{article}
\usepackage[T1]{fontenc}
\usepackage[utf8]{inputenc}
\usepackage[spanish]{babel}
\usepackage{geometry}
\usepackage{graphicx}
\usepackage{xcolor}
\usepackage{listings}
\usepackage{booktabs}
\usepackage{hyperref}
\usepackage{palatino}
\usepackage{amsmath}
\usepackage{textcomp}

\geometry{top=2.5cm, bottom=2.5cm, left=2.5cm, right=2.5cm}

% Configuración de colores y estilo para el código
\definecolor{codegreen}{rgb}{0,0.6,0}
\definecolor{codegray}{rgb}{0.5,0.5,0.5}
\definecolor{codepurple}{rgb}{0.58,0,0.82}
\definecolor{backcolour}{rgb}{0.95,0.95,0.92}

\lstdefinestyle{mystyle}{
    backgroundcolor=\color{backcolour},
    commentstyle=\color{codegreen},
    keywordstyle=\color{magenta},
    numberstyle=\tiny\color{codegray},
    stringstyle=\color{codepurple},
    basicstyle=\ttfamily\footnotesize,
    breakatwhitespace=false,
    breaklines=true,
    captionpos=b,
    keepspaces=true,
    numbers=left,
    numbersep=5pt,
    showspaces=false,
    showstringspaces=false,
    showtabs=false,
    tabsize=2
}

\lstset{style=mystyle,
    literate={á}{{\'a}}1
             {é}{{\'e}}1
             {í}{{\'i}}1
             {ó}{{\'o}}1
             {ú}{{\'u}}1
             {ñ}{{\~n}}1
             {Á}{{\'A}}1
             {É}{{\'E}}1
             {Í}{{\'I}}1
             {Ó}{{\'O}}1
             {Ú}{{\'U}}1
             {Ñ}{{\~N}}1
             {¡}{{\textexclamdown}}1
}

\title{\textbf{Flujo de Trabajo para Fabricación de PCBs}\\ \large De KiCad a G-Code con Python y CNC}
\author{Documentación generada por Gemini}
\date{\today}

\begin{document}

\maketitle
\tableofcontents
\newpage

\section{Del Esquema al G-Code con Python}

Es totalmente posible generar G-Code para una CNC a partir de un diagrama esquemático utilizando librerías de Python. Sin embargo, es un camino que requiere un par de pasos intermedios, ya que las máquinas CNC no pueden "ver" un esquema eléctrico (que es una representación lógica) y saber dónde van las pistas físicas.

El flujo técnico es: \textbf{Esquema $\rightarrow$ Layout de PCB (Gerber) $\rightarrow$ G-Code}.

\subsection{La pieza clave: El formato Gerber}

Antes de generar G-Code, necesitas que tu esquema se convierta en un diseño de pistas físicas. En la industria, este estándar se llama \textbf{Gerber (.gbr)}. Python es excelente para manipular estos archivos.

\subsection{Librerías de Python recomendadas}

\begin{itemize}
    \item \textbf{Para procesar archivos Gerber: \texttt{pcb-tools-extension}}. Esta es la librería estándar para leer y manipular archivos de fabricación de circuitos. Te permite extraer las coordenadas de las pistas y los pads.
    \item \textbf{Para generar el G-Code: \texttt{gcodepy}}. Es una librería ligera para escribir líneas de G-Code de forma programática (especificar velocidades de avance, levantamiento de la broca y fresado).
    \item \textbf{El "Todo en uno": \texttt{pcb2gcode}}. Es una herramienta de línea de comandos (escrita en C++) que se puede invocar desde Python. Es el estándar de oro para convertir Gerbers a G-Code de aislamiento.
\end{itemize}

\section{El Flujo Maestro: De la Idea al Cobre}

El proceso completo, uniendo el diseño en KiCad con la automatización en Python, es el siguiente:

\begin{enumerate}
    \item \textbf{KiCad (Diseño):} Dibujas tu esquema y luego haces el PCB Layout (donde ubicas físicamente los componentes y trazas las rutas).
    \item \textbf{Exportación (Gerber):} Generas los archivos \texttt{.gbr}. Para una CNC básica, normalmente solo necesitas:
    \begin{itemize}
        \item \texttt{F\_Cu.gbr} (Capa de cobre frontal) o \texttt{B\_Cu.gbr} (Posterior).
        \item \texttt{Edge\_Cuts.gbr} (El borde de la placa).
        \item \texttt{-drl.gbr} (Archivo de taladrado para los agujeros).
    \end{itemize}
    \item \textbf{Procesado con Python:} Aquí es donde entra tu script. La librería lee las coordenadas geométricas del Gerber.
    \item \textbf{Generación de G-Code:} Tu script traduce esas geometrías en comandos \texttt{G0} (movimiento rápido) y \texttt{G1} (corte).
\end{enumerate}

\subsection{El gran "Truco" técnico: El Aislamiento (Isolation Routing)}

Este es el punto donde la mayoría falla al principio. Los archivos Gerber definen dónde \textbf{hay} cobre. Sin embargo, la CNC necesita saber por dónde pasar la fresa para \textbf{separar} el cobre. Si tu pista mide 0.5mm y tu fresa mide 0.2mm, no puedes pasar la fresa por el centro de la pista, ¡porque la destruirías! Tu script de Python debe calcular un \textit{offset} (desplazamiento).

\section{Auto-enrutado (Autorouting) con la API de Python de KiCad}

KiCad tiene una librería interna de Python llamada \texttt{pcbnew}. No necesitas exportar nada; el script corre dentro de tu archivo de proyecto \texttt{.kicad\_pcb}. Con ella puedes leer la "Netlist", mover componentes y trazar pistas automáticamente.

\subsection{Script de Ejemplo: Creando una pista}

Este script dibuja una pista (track) uniendo dos puntos específicos en la capa de cobre superior (\texttt{F\_Cu}).

\begin{lstlisting}[language=Python, caption=Creando una pista programáticamente en KiCad]
import pcbnew

# 1. Obtener el diseno (Board) actual que tienes abierto
board = pcbnew.GetBoard()

# 2. Definir los puntos de inicio y fin (en milimetros)
start_point = pcbnew.VECTOR2I_MM(100, 100)
end_point = pcbnew.VECTOR2I_MM(120, 110)

# 3. Crear el objeto de la pista (Track)
track = pcbnew.PCB_TRACK(board)

# 4. Configurar las propiedades de la pista
track.SetStart(start_point)
track.SetEnd(end_point)
track.SetWidth(pcbnew.FromMM(0.5)) # Ancho de 0.5 mm
track.SetLayer(pcbnew.F_Cu)        # Capa de cobre frontal

# 5. Anadir la pista a la placa
board.Add(track)

# 6. Refrescar la interfaz para ver los cambios
pcbnew.Refresh()

print("¡Pista creada con éxito!")
\end{lstlisting}

\section{Implementando un Auto-router Básico}

\subsection{Escaneo de Componentes y Pads}

Para que un script sea "inteligente", primero debe saber quién es quién en la placa. Este código extrae la "hoja de ruta" de tu PCB.

\begin{lstlisting}[language=Python, caption=Mapeo de componentes y sus terminales]
import pcbnew

def mapear_componentes():
    board = pcbnew.GetBoard()
    print(f"{'Ref':<8} | {'Pad':<5} | {'X (mm)':<10} | {'Y (mm)':<10} | {'Red (Net)'}")
    print("-" * 60)

    for footprint in board.GetFootprints():
        ref = footprint.GetReference() # Ejemplo: 'R1', 'U1'
        for pad in footprint.Pads():
            num_pad = pad.GetNumber()
            pos = pad.GetPosition() # Coordenadas en nanometros
            
            x_mm = pcbnew.ToMM(pos.x)
            y_mm = pcbnew.ToMM(pos.y)
            net_name = pad.GetNetname()
            
            print(f"{ref:<8} | {num_pad:<5} | {x_mm:<10.3f} | {y_mm:<10.3f} | {net_name}")

mapear_componentes()
\end{lstlisting}

\subsection{Creando una Cuadrícula de Obstáculos (Discretización)}

Para que un script pueda "pensar" por dónde pasar una pista, convertimos el mundo continuo de milímetros en un mapa de celdas (una rejilla o \textit{grid}). Cada cuadro puede estar "Libre" (0) o "Ocupado" (1).

\begin{lstlisting}[language=Python, caption=Creación de la matriz de obstáculos]
import pcbnew

RESOLUCION = 0.25 # mm por celda
ANCHO_PLACA = 50  # mm
ALTO_PLACA = 50   # mm

filas = int(ALTO_PLACA / RESOLUCION)
columnas = int(ANCHO_PLACA / RESOLUCION)
grid = [[0 for _ in range(columnas)] for _ in range(filas)]

def marcar_obstaculos(board):
    for footprint in board.GetFootprints():
        for pad in footprint.Pads():
            pos = pad.GetPosition()
            x_mm = pcbnew.ToMM(pos.x)
            y_mm = pcbnew.ToMM(pos.y)
            
            grid_x = int(x_mm / RESOLUCION)
            grid_y = int(y_mm / RESOLUCION)
            
            if 0 <= grid_x < columnas and 0 <= grid_y < filas:
                grid[grid_y][grid_x] = 1 # Marcar como obstáculo

board = pcbnew.GetBoard()
marcar_obstaculos(board)
\end{lstlisting}

\subsection{Búsqueda de Caminos con A* (A-Estrella)}

Ahora que tienes el grid, necesitas un algoritmo que encuentre el camino del Punto A al Punto B evitando los obstáculos. El algoritmo \textbf{A* (A-Estrella)} es el estándar de oro. Funciona evaluando un costo: $f(n) = g(n) + h(n)$.

\begin{lstlisting}[language=Python, caption=Búsqueda de camino con la librería `pathfinding`]
from pathfinding.core.diagonal_movement import DiagonalMovement
from pathfinding.core.grid import Grid
from pathfinding.finder.a_star import AStarFinder

# 1. Crear el mapa de la placa (Matriz de pesos)
# 1 = Transitable, 0 = Obstáculo
matrix = [
    [1, 1, 1, 1, 1],
    [1, 0, 0, 0, 1], # Barrera de obstáculos
    [1, 1, 1, 0, 1],
    [1, 1, 1, 1, 1]
]

# 2. Inicializar el grid de la librería
grid = Grid(matrix=matrix)

# 3. Definir puntos de inicio y fin
start = grid.node(0, 0)
end = grid.node(4, 3)

# 4. Configurar y ejecutar el buscador (A*)
finder = AStarFinder(diagonal_movement=DiagonalMovement.always)
path, runs = finder.find_path(start, end, grid)

print(f"Ruta encontrada: {path}")
\end{lstlisting}

\section{Generación del G-Code Final}

Esta es la pieza final del rompecabezas: convertir esa lista de coordenadas abstractas en un archivo físico que tu máquina CNC pueda interpretar.

\begin{lstlisting}[language=Python, caption=Generador de G-Code a partir de una ruta]
def generar_gcode(path, resolucion, archivo_salida="placa.nc"):
    Z_SEGURIDAD = 2.0
    Z_CORTE = -0.1
    F_AVANCE = 300

    with open(archivo_salida, "w") as f:
        f.write("G21 ; Unidades en milimetros\n")
        f.write("G90 ; Posicionamiento absoluto\n")
        f.write(f"G0 Z{Z_SEGURIDAD} ; Subir fresa\n\n")

        if not path: return

        inicio_x = path[0][0] * resolucion
        inicio_y = path[0][1] * resolucion
        f.write(f"G0 X{inicio_x:.3f} Y{inicio_y:.3f}\n")
        f.write(f"G1 Z{Z_CORTE} F100\n")

        for i in range(1, len(path)):
            x = path[i][0] * resolucion
            y = path[i][1] * resolucion
            f.write(f"G1 X{x:.3f} Y{y:.3f} F{F_AVANCE}\n")

        f.write(f"\nG0 Z{Z_SEGURIDAD}\n")
        f.write("M2 ; Fin del programa\n")

mi_ruta = [(0,0), (1,1), (2,1), (3,2)]
generar_gcode(mi_ruta, resolucion=0.25)
\end{lstlisting}

\section{Técnica Avanzada: Auto-Nivelación (Auto-Leveling)}

El Auto-leveling es la diferencia entre un circuito profesional y uno con pistas rotas. Como el cobre de una PCB es muy delgado (0.035 mm), cualquier curvatura arruina el trabajo. La solución es crear un \textbf{Mapa de Alturas}.

\begin{enumerate}
    \item \textbf{Muestreo (Probing):} Antes de fresar, el script le indica a la CNC que toque la placa en varios puntos de una cuadrícula. La máquina usa la broca y la placa como un interruptor para registrar la altura $Z$ exacta en cada punto $(X, Y)$. Esto se hace con el comando \texttt{G38.2}.
    \item \textbf{Compensación (Interpolación):} Cuando se va a fresar un punto, Python calcula la altura exacta en ese lugar usando una interpolación bilineal, promediando las alturas de los 4 puntos de la cuadrícula más cercanos. La altura final se ajusta: $Z_{adj} = Z_{original} + Z_{offset}(x, y)$.
\end{enumerate}

\section{Comunicación con la CNC: El G-Code Sender}

La mayoría de las CNCs de código abierto utilizan \textbf{GRBL}. Este firmware escucha a través de un puerto serie (USB) y espera que le envíes una línea de G-Code, la procese y responda con un "ok" para recibir la siguiente.

\begin{lstlisting}[language=Python, caption=Python G-Code Sender para GRBL]
import serial
import time

def enviar_gcode(puerto_com, archivo_nc):
    try:
        s = serial.Serial(puerto_com, 115200)
        print(f"Conectado a la CNC en {puerto_com}")
    except Exception as e:
        print(f"Error al conectar: {e}")
        return

    s.write(b"\r\n\r\n") # Despertar a GRBL
    time.sleep(2)
    s.flushInput()

    with open(archivo_nc, 'r') as f:
        for linea in f:
            l = linea.strip()
            if l:
                print(f"Enviando: {l}")
                s.write((l + '\n').encode('utf-8'))
                
                # Esperar respuesta "ok" de la CNC
                respuesta = s.readline()
                print(f"Respuesta: {respuesta.decode('utf-8').strip()}")

    s.close()
    print("Trabajo finalizado.")

# Uso:
# enviar_gcode('/dev/ttyUSB0', 'placa.nc') # En Linux
\end{lstlisting}

\section{Conclusión: Soberanía Tecnológica con Software Libre}

El hecho de que se pueda pasar de un concepto abstracto de ingeniería a una pieza física utilizando un stack de software 100\% libre (\textbf{KiCad + Python + GRBL}) es sumamente potente. No se depende de licencias costosas ni de formatos cerrados.

Este flujo de trabajo te da control total sobre el algoritmo, desde la geometría de aislamiento hasta la gestión de residuos, abriendo la puerta a la verdadera soberanía tecnológica.

\end{document}